% ============================================================================
% GLOSA-7SSA GLOBAL SCHOLARLY PUBLICATION SCHEMA
% Standalone LaTeX master template
% Version 1.0 | 2026-09-25
%
% PURPOSE
% -------
% A journal-facing + audit-facing manuscript architecture that integrates:
%   (1) glosa's claim/evidence/provenance discipline,
%   (2) 7SSA: Seven-Section Scholarly Architecture,
%   (3) leading international journal declaration practices,
%   (4) Thai NRCT guidance on ethical use of Generative AI in research,
%   (5) cross-publisher AI, authorship, COI, funding, ethics, data/code,
%       reproducibility and transparency statements.
%
% DESIGN RULE
% -----------
% Visible journal sections may change, but the seven intellectual functions
% must remain. Internal glosa audit material is switchable and should normally
% be submitted only when the target journal requests or permits it.
%
% POLICY NOTE
% -----------
% This is a conservative cross-publisher interoperability profile, NOT a claim
% that one single universal publisher standard exists. The current instructions
% of the target journal always override placement, wording, file-format, and
% disclosure details.
%
% COMPILATION
% -----------
% Tested as a self-contained pdflatex document. For visible Thai script in a
% journal-specific derivative, use XeLaTeX/LuaLaTeX and the journal's permitted
% font setup. Thai text is therefore kept mainly in comments in this portable
% pdflatex master.
%
% SOURCE ANCHORS USED IN THE DESIGN
% ---------------------------------
% - morrocwi/glosa current public repository:
%   methodology/P07_ai_fill_disclosure.md
%   methodology/P10_publish_gate.md
%   methodology/P13_literature_review.md
%   templates/paper/arxiv-onecol/main.tex
%   templates/paper/arxiv-twocol/main.tex
% - NRCT (Thailand), Ethical Application of Generative AI for Researchers,
%   September 2026.
% - ICMJE Recommendations, updated January 2026.
% - Elsevier Generative AI policies for journals, updated 2026.
% - Springer Nature AI editorial/manuscript policies, current 2026.
% - Wiley Best Practice Guidelines on Research Integrity and Publishing Ethics.
% - ANSI/NISO CRediT taxonomy (14 contributor roles).
%
% ============================================================================

\documentclass[11pt]{article}

% ---- Portable packages ------------------------------------------------------
\usepackage[margin=1in]{geometry}
\usepackage[T1]{fontenc}
\usepackage[utf8]{inputenc}
\usepackage{lmodern}
\usepackage{microtype}
\usepackage{amsmath,amssymb,mathtools}
\usepackage{booktabs}
\usepackage{tabularx}
\usepackage{array}
\usepackage{longtable}
\usepackage{enumitem}
\usepackage{graphicx}
\usepackage{xcolor}
\usepackage{url}
\usepackage{hyperref}
\usepackage{cleveref}
\usepackage{titlesec}
\usepackage{caption}
\usepackage{authblk}
\usepackage{fancyhdr}
\usepackage{float}
\usepackage{ifthen}
\usepackage{xparse}

% ---- Appearance -------------------------------------------------------------
\definecolor{Ink}{RGB}{25,25,28}
\definecolor{Slate}{RGB}{82,88,96}
\definecolor{Soft}{RGB}{245,245,245}
\hypersetup{
  colorlinks=true,
  linkcolor=Ink,
  citecolor=Ink,
  urlcolor=Slate,
  pdftitle={GLOSA-7SSA Global Scholarly Publication Schema},
  pdfauthor={[FILL: Human author(s)]}
}
\setlength{\parindent}{1em}
\setlength{\parskip}{0.25em}
\setlist{nosep,leftmargin=*}
\captionsetup{font=small,labelfont=bf}
\titleformat{\section}{\large\bfseries}{\thesection}{0.6em}{}
\titleformat{\subsection}{\normalsize\bfseries}{\thesubsection}{0.5em}{}

% ---- Modes -----------------------------------------------------------------
% submissionmode:
%   true  = journal-facing manuscript
%   false = internal development/audit copy
%
% glosaaudit:
%   true  = append glosa claim matrix, AI-route audit, mastery gate, etc.
%   false = omit internal audit appendices
%
% thaiarticle:
%   true  = semantic flag for Thai-journal renderer; headings remain English in
%           this pdflatex master for portability. Replace heading strings in a
%           XeLaTeX/LuaLaTeX derivative if Thai script is required.
\newif\ifsubmissionmode
\newif\ifglosaaudit
\newif\ifthaiarticle

\submissionmodetrue
\glosaaudittrue
\thaiarticlefalse

% ---- Schema metadata --------------------------------------------------------
\newcommand{\SchemaName}{GLOSA--7SSA Global Scholarly Publication Schema}
\newcommand{\SchemaVersion}{1.0}
\newcommand{\SchemaDate}{2026-09-25}
\newcommand{\TargetVenue}{[FILL: journal / venue]}
\newcommand{\ArticleType}{[FILL: conceptual / theory / review / legal / formal / SoK / policy / other]}
\newcommand{\TargetSystem}{[FILL: Q1 World / Thai leading journal / repository working paper]}
\newcommand{\PolicyCheckedDate}{[FILL: YYYY-MM-DD]}

% ---- Glosa tier macros -------------------------------------------------------
\newcommand{\ThCoqc}{\textsc{Th\_coqc}}
\newcommand{\FiniteDiag}{\textsc{finite\_diagnostic}}
\newcommand{\FitCal}{\textsc{fit\_calibrated}}
\newcommand{\DrTier}{\textsc{Dr}}
\newcommand{\DefTier}{\textsc{definition}}
\newcommand{\OpenTier}{\textsc{Open}}
\newcommand{\tier}[1]{\textsc{#1}}
\newcommand{\claimref}[1]{\textsuperscript{[#1]}}

% ---- Status labels -----------------------------------------------------------
\newcommand{\Required}{\textbf{REQUIRED}}
\newcommand{\Conditional}{\textbf{CONDITIONAL}}
\newcommand{\Internal}{\textbf{INTERNAL}}
\newcommand{\NA}{\textit{Not applicable}}
\newcommand{\Fill}{\texttt{[FILL]}}
\newcommand{\Pass}{\textsc{PASS}}
\newcommand{\Hold}{\textsc{HOLD}}
\newcommand{\Fail}{\textsc{FAIL}}

% ---- Utility environments ---------------------------------------------------
\newenvironment{schemaBox}[1]
  {\begin{center}\begin{minipage}{0.96\linewidth}\small
   \noindent\textbf{#1}\par\medskip}
  {\end{minipage}\end{center}}

\newenvironment{declarationblock}[1]
  {\subsection*{#1}\addcontentsline{toc}{subsection}{#1}}
  {}

\newenvironment{auditblock}[1]
  {\subsection*{#1}\addcontentsline{toc}{subsection}{#1}\small}
  {\normalsize}

% ---- Header/footer -----------------------------------------------------------
\pagestyle{fancy}
\fancyhf{}
\fancyhead[L]{\footnotesize \SchemaName}
\fancyhead[R]{\footnotesize v\SchemaVersion\ --- \ArticleType}
\fancyfoot[C]{\footnotesize \thepage}
\renewcommand{\headrulewidth}{0.3pt}

% ============================================================================
% FRONT MATTER — PUBLICATION IDENTITY LAYER
% ============================================================================

\title{\textbf{[FILL: Manuscript Title]}\\[0.35em]
\large [FILL: Optional subtitle]}
\author[1]{[FILL: Author One]\thanks{ORCID: [FILL]}}
\author[2]{[FILL: Author Two]\thanks{ORCID: [FILL]}}
\affil[1]{[FILL: Affiliation or ``Independent Researcher''] }
\affil[2]{[FILL: Affiliation]}
\date{}

\begin{document}
\maketitle

\begin{schemaBox}{Submission profile}
\begin{tabularx}{\linewidth}{@{}lX@{}}
\toprule
Schema & \SchemaName\ v\SchemaVersion\ (\SchemaDate)\\
Target system & \TargetSystem\\
Target venue & \TargetVenue\\
Article type & \ArticleType\\
Journal policy last checked & \PolicyCheckedDate\\
Corresponding author & [FILL: name, email]\\
Manuscript status & [FILL: draft / preprint / submitted / revision]\\
Registration / protocol / preregistration & [FILL: DOI/URL/id or N/A]\\
Repository / archival record & [FILL: DOI/URL or N/A]\\
\bottomrule
\end{tabularx}
\end{schemaBox}

\begin{abstract}
[FILL: 150--250 words, adapted to the target journal.

Functional order:
(1) important problem;
(2) exact limitation/gap;
(3) approach or scholarly method;
(4) principal new contribution;
(5) main implication and one important boundary/limitation.

Do not claim evidence, novelty, causality, validation, or generality beyond the
body of the manuscript.]
\end{abstract}

\noindent\textbf{Keywords:} [FILL: 4--8 terms]

% Optional publisher items: highlights, graphical abstract, classifications.
% Keep only if the target journal asks for them.

% ============================================================================
% THE 7SSA CORE — SEVEN INTELLECTUAL SECTORS
% ============================================================================


\section*{Core Epistemic Structure}
\label{sec:ces}
Per \texttt{methodology/P20\_core\_epistemic\_structure.md}: who brought the world in, who bridged,
and which machines took part --- three roles, never merged (non-collapse rule below). Registration
is an attribution device, not a certification of truth or expertise.
\begin{description}
  \item[Core Respondent / Experience-Based Expert:] [FILL: name or anonymised role] ---
    [FILL: lived experience; selections; interpretation]
  \item[Interactional Expert:] [FILL: name / role, or \textbf{None}]
  \item[AI Model(s) Used:] [FILL: Model 1 --- role]; [FILL: Model 2 --- role] (role disclosure only; never authorship)
\end{description}
Non-collapse rule: Experience-Based Expertise $\neq$ Interactional Expertise $\neq$ AI Model.

\section{Introduction}
\label{sec:s1}

% THAI RENDERER:
% 1. บทนำ
%
% Intellectual function:
% Problem + importance + scholarly conversation + gap preview + central argument
% + exact contribution.

\subsection{Problem and phenomenon}
[FILL: State the phenomenon or intellectual problem without smuggling the
proposed answer into the problem statement.]

\subsection{Why the problem matters}
[FILL: Explain why the target field or readership should care. Separate
disciplinary importance from local/practical importance.]

\subsection{Scholarly conversation}
[FILL: Locate the active conversation, including the closest international
and, where relevant, Thai/local scholarship.]

\subsection{Unresolved tension}
[FILL: Preview the contradiction, missing mechanism, conceptual ambiguity,
boundary failure, assumption failure, new phenomenon, or other scholarly gap.]

\subsection{Central argument and contribution}
\noindent\textbf{Central argument.}
[FILL: one sentence beginning, if useful, ``We argue that ...'']

\noindent\textbf{Contribution.}
[FILL: name the new intellectual object: construct, distinction, theory,
framework, mechanism, taxonomy, ontology, architecture, formalization,
interpretation, legal principle, or research agenda.]

\noindent\textbf{Roadmap.}
[FILL: optional; follow target-journal style.]

% ---------------------------------------------------------------------------

\section{Approach, Scope, and Knowledge Base}
\label{sec:s2}

% THAI RENDERER:
% 2. แนวทาง ขอบเขต และฐานความรู้
%
% This is not necessarily "Methods" in a conceptual paper. It is the
% reconstruction layer: how the paper knows what it claims to know.

\subsection{Scholarly method}
[FILL: conceptual analysis / theory construction / philosophical argument /
integrative synthesis / doctrinal legal analysis / formal method /
systematization of knowledge / critical review / policy analysis / other.]

\subsection{Scope and exclusions}
\begin{itemize}
  \item Domain/population/object: [FILL]
  \item Geography/jurisdiction: [FILL]
  \item Time window: [FILL]
  \item Languages/source families: [FILL]
  \item Explicit exclusions: [FILL]
\end{itemize}

\subsection{Knowledge-base construction}
[FILL: databases/search routes, canonical sources, target-journal literature,
adjacent disciplines, local literature, source-acquisition logic, and how
sources were verified.]

\subsection{Literature-review mode}
\label{sec:lrs-mode}
Select honestly:
\begin{itemize}
  \item [FILL] Targeted scholarly search
  \item [FILL] Integrative review
  \item [FILL] Systematic review / scoping review / SoK (only if protocol supports it)
  \item [FILL] Doctrinal/legal corpus
  \item [FILL] Archival/historical corpus
\end{itemize}

\subsection{AI involvement in the research process}
\label{sec:ai-method}
If AI affected the research process (not merely manuscript language), describe
here the tool/model, version or release identifier when available, provider,
date(s) of use, research stage, inputs, outputs, validation, and retained logs.
The public declaration in \Cref{sec:ai-declaration} summarizes manuscript-level
AI use; this subsection gives method-level reproducibility detail.

\begin{schemaBox}{Glosa compatibility rule}
AI-generated content is never treated as evidence merely because the model
produced it. AI-discovered or AI-suggested claims, sources, analyses, code, or
interpretations enter the manuscript only after the relevant human/source/
formal/computational verification route has been completed and logged.
\end{schemaBox}

% ---------------------------------------------------------------------------

\section{Existing Knowledge and Theoretical Foundation}
\label{sec:s3}

% THAI RENDERER:
% 3. แนวคิด ทฤษฎี และองค์ความรู้ที่เกี่ยวข้อง

\subsection{Key constructs}
[FILL: Define the constructs used by the argument; distinguish neighboring
constructs and levels of analysis.]

\subsection{Major theories or positions}
[FILL: Organize by problems/positions/mechanisms, not by a chronological
``Author A said; Author B said'' list.]

\subsection{Nearest intellectual neighbors}
[FILL: Identify the strongest prior theory/framework/method that could make
this paper unnecessary.]

\subsection{What is already established}
[FILL: State K0---the knowledge state before this paper---with claim-matched
citations.]

% ---------------------------------------------------------------------------

\section{Problem in Existing Knowledge}
\label{sec:s4}

% THAI RENDERER:
% 4. ปัญหา ข้อจำกัด หรือช่องว่างขององค์ความรู้เดิม

\subsection{Exact inadequacy}
Select and justify one or more:
\begin{itemize}
  \item conceptual ambiguity;
  \item contradiction;
  \item missing mechanism;
  \item fragmented literatures;
  \item unclear boundary conditions;
  \item hidden/failed assumption;
  \item formalization gap;
  \item classification gap;
  \item level-of-analysis gap;
  \item normative or doctrinal gap;
  \item new phenomenon not handled by existing theory.
\end{itemize}

\subsection{Why the inadequacy matters}
[FILL: What can the field not explain, distinguish, justify, predict, organize,
or decide because of this problem?]

\subsection{Strongest prior attempt}
[FILL: Treat the strongest existing solution charitably. Explain why it is
still insufficient without manufacturing a gap.]

% ---------------------------------------------------------------------------

\section{New Contribution}
\label{sec:s5}

% THAI RENDERER:
% 5. การวิเคราะห์และข้อเสนอใหม่ของผู้เขียน

\subsection{Name and definition of the contribution}
\noindent\textbf{Name:} [FILL]\\
\textbf{Type:} [FILL: construct / distinction / mechanism / theory / framework /
taxonomy / ontology / architecture / principle / formalization / interpretation]\\
\textbf{Definition:} [FILL]

\subsection{Components and relations}
[FILL: Explain nodes, entities, dimensions, mechanisms, transitions, or
inferential steps.]

\subsection{Derivation}
[FILL: Show how the contribution follows from Sectors 3 and 4 rather than
appearing as an unsupported diagram.]

\subsection{Propositions, principles, or formal statements}
\begin{enumerate}
  \item P1: [FILL]
  \item P2: [FILL]
  \item P3: [FILL]
\end{enumerate}

\subsection{Comparison with nearest prior framework}
[FILL: same / different / extends / challenges. Avoid unsupported priority
claims such as ``first ever'' unless an adequate search supports them.]

% ---------------------------------------------------------------------------

\section{Critical Evaluation, Boundaries, and Implications}
\label{sec:s6}

% THAI RENDERER:
% 6. การอภิปราย ข้อโต้แย้ง ข้อจำกัด และนัยสำคัญ

\subsection{Strongest objection}
[FILL: State the strongest plausible objection in a form its proponent would
recognize.]

\subsection{Alternative explanation or rival theory}
[FILL]

\subsection{Counterexample or failure case}
[FILL]

\subsection{Response and repair}
[FILL: Defend, revise, or narrow the contribution. A narrow honest theory is
better than an unfalsifiable broad one.]

\subsection{Boundary conditions}
[FILL: Where should the contribution apply? Where should it not apply?]

\subsection{Limitations}
[FILL: epistemic, conceptual, methodological, empirical, legal, cultural,
technical, data, AI, and generalizability limits as applicable.]

\subsection{Implications}
\begin{itemize}
  \item Theoretical implications: [FILL]
  \item Methodological implications: [FILL]
  \item Practical/professional implications: [FILL or N/A]
  \item Policy/governance implications: [FILL or N/A]
  \item Future research: [FILL]
\end{itemize}

% ---------------------------------------------------------------------------

\section{Conclusion}
\label{sec:s7}

% THAI RENDERER:
% 7. บทสรุป

\subsection*{Answer to the original problem}
[FILL]

\subsection*{Before $\rightarrow$ After}
\begin{description}[style=nextline,leftmargin=1.4em]
  \item[Before this paper:] [FILL: what the field could not adequately
  distinguish/explain/organize.]
  \item[After this paper:] [FILL: the precise new capability supplied by the
  contribution.]
\end{description}

\subsection*{What this paper does not claim}
[FILL: explicit non-claims.]

\subsection*{Next frontier}
[FILL]

% ============================================================================
% WORLD-STANDARD DECLARATIONS & TRANSPARENCY LAYER
% These are NOT "Sector 8". They are publication-governance metadata around
% the 7SSA intellectual core.
% ============================================================================

\section*{Statements and Declarations}
\addcontentsline{toc}{section}{Statements and Declarations}

\begin{declarationblock}{Author contributions (CRediT)}
% Use the target journal's exact implementation. CRediT has 14 roles.
% CRediT describes contributions; it does not itself decide who qualifies
% for authorship.
\begin{tabularx}{\linewidth}{@{}p{0.31\linewidth}X@{}}
\toprule
CRediT role & Contributor(s)\\
\midrule
Conceptualization & [FILL]\\
Data curation & [FILL / N/A]\\
Formal analysis & [FILL / N/A]\\
Funding acquisition & [FILL / N/A]\\
Investigation & [FILL / N/A]\\
Methodology & [FILL]\\
Project administration & [FILL / N/A]\\
Resources & [FILL / N/A]\\
Software & [FILL / N/A]\\
Supervision & [FILL / N/A]\\
Validation & [FILL]\\
Visualization & [FILL / N/A]\\
Writing -- original draft & [FILL]\\
Writing -- review \& editing & [FILL]\\
\bottomrule
\end{tabularx}
\end{declarationblock}

\begin{declarationblock}{Authorship and accountability}
[FILL: Confirm that every listed author satisfies the target journal's
authorship criteria, approved the submitted version, and accepts responsibility
for the work. AI systems/tools are not listed as authors.]
\end{declarationblock}

\begin{declarationblock}{Declaration of generative AI and AI-assisted technologies}
\label{sec:ai-declaration}

% Conservative cross-publisher superset.
% At submission, adapt placement and exact wording to the target journal.
%
% For glosa: keep the detailed route-level record internally even if the
% publisher allows a shorter public declaration.

\noindent\textbf{Public summary.}
During the preparation and/or research process for this work, the author(s)
used [FILL: tool/service/model] for [FILL: purpose and research/manuscript
stage]. The author(s) reviewed, verified, and edited all material as required
and retain full responsibility for the submitted work.

\medskip
\noindent\textbf{Detailed disclosure record.}

\begin{longtable}{@{}p{0.28\linewidth}p{0.66\linewidth}@{}}
\toprule
Field & Disclosure\\
\midrule
\endhead
Tool/service & [FILL: product/service/model]\\
Provider/developer & [FILL]\\
Model/version/release & [FILL: where available]\\
Dates accessed & [FILL]\\
Use category & [FILL: ideation; information discovery; refinement; translation/editing;
analysis/recommendation; code generation/optimization; visual/media creation; other]\\
Research/manuscript stage & [FILL: question; literature; design; code; data; analysis;
visualization; drafting; review; submission preparation]\\
Purpose & [FILL]\\
Extent/influence & [FILL: did it affect central arguments, methods, results, or conclusions?]\\
Input-data class & [FILL: public / licensed / internal / confidential / personal /
sensitive / unpublished / copyrighted / none]\\
Prompt/settings retained & [FILL: yes/no; location; explain any lawful/privacy redactions]\\
Outputs retained & [FILL: yes/no; location]\\
Human verification & [FILL: what was checked; against what independent source/method]\\
Human checker/approver & [FILL]\\
Citation verification & [FILL: source-first verification route]\\
Code verification & [FILL: tests/oracle/review, or N/A]\\
Visual-content verification & [FILL: underlying verifiable source/data, or N/A]\\
Privacy/IP/terms check & [FILL]\\
Material effect on conclusion & [FILL: none / limited / substantive; explain]\\
\bottomrule
\end{longtable}

\noindent\textbf{Non-delegation statement.}
Human author(s) retain responsibility for the scholarly argument, interpretation,
source verification, authorship decisions, and final submission. AI output is
not used as a primary scholarly source and is not treated as evidence merely
because it was generated.
\end{declarationblock}

\begin{declarationblock}{Funding}
[FILL one:
(a) This work received support from [funder, grant number], and the funder had
[role/no role] in [design/analysis/writing/decision to submit];
or
(b) The authors received no specific funding for this work.]
\end{declarationblock}

\begin{declarationblock}{Competing interests / conflict of interest}
% Disclose relationships or activities that could reasonably be perceived to
% influence the work. Follow the target journal's disclosure period and form.
%
% Consider: financial interests, employment, consultancy, advisory roles,
% speaking fees, stocks/shares, patents/IP, litigation, leadership roles,
% personal/professional relationships, political/religious/advocacy interests
% where directly relevant, or other non-financial competing interests.
[FILL: Declare relevant competing interests, or state that none are declared.]
\end{declarationblock}

\begin{declarationblock}{Ethics approval and research integrity}
[FILL according to the work:
- approving REC/IRB/ethics body and approval number;
- waiver/exemption and basis;
- for a purely conceptual/review paper with no new human/animal data, state
  the applicable N/A rationale;
- for community/indigenous/culturally sensitive knowledge, identify relevant
  consent/governance route where applicable.]
\end{declarationblock}

\begin{declarationblock}{Consent to participate}
[FILL: approval/consent details or N/A.]
\end{declarationblock}

\begin{declarationblock}{Consent for publication}
[FILL: consent for identifiable material, quotations, images, case details, or N/A.]
\end{declarationblock}

\begin{declarationblock}{Data availability}
[FILL one:
- data DOI/repository and access conditions;
- restricted data + justified access route;
- no new data were created or analyzed;
- all data are contained within article/supplement.]
\end{declarationblock}

\begin{declarationblock}{Code and software availability}
[FILL: repository, commit/tag, license, environment/lock file, DOI, or N/A.]
\end{declarationblock}

\begin{declarationblock}{Materials, protocol, and preregistration}
[FILL: protocol DOI/registration, materials repository, preregistration id,
or N/A. Do not imply preregistration if the protocol was frozen after results
were seen.]
\end{declarationblock}

\begin{declarationblock}{Reproducibility and verification}
[FILL: identify what another scholar can inspect or rerun: search manifest,
source ledger, theorem proof, software environment, data pipeline, analysis
script, simulation, benchmark, claim-evidence matrix, etc.]
\end{declarationblock}

\begin{declarationblock}{Intellectual property, permissions, and third-party material}
[FILL: permissions/licenses for copyrighted figures/tables/text/data; patent
or confidential-information considerations; or N/A.]
\end{declarationblock}

\begin{declarationblock}{Dual-use, safety, and biosecurity statement}
[FILL where relevant: safety review, controlled details, risk mitigation, or N/A.]
\end{declarationblock}

\begin{declarationblock}{Acknowledgements}
[FILL: people, organizations, facilities, language/editorial assistance,
non-author contributors, and permissions to acknowledge where required.]
\end{declarationblock}

\begin{declarationblock}{Prior dissemination, preprint, and overlapping publication}
[FILL: preprint/working-paper DOI, conference abstract, thesis chapter,
translation, prior version, or state none. Explain material overlap.]
\end{declarationblock}

\begin{declarationblock}{Reporting guideline}
[FILL if applicable: PRISMA, CONSORT, STROBE, CARE, SRQR, COREQ, ARRIVE,
TRIPOD, CHEERS, etc.; otherwise N/A.]
\end{declarationblock}

% ============================================================================
% REFERENCES
% Replace this placeholder with journal-specific bibliography handling.
% ============================================================================

\section*{References}
\addcontentsline{toc}{section}{References}

\begin{thebibliography}{9}

\bibitem{placeholder1}
[FILL: Reference 1.]

\bibitem{placeholder2}
[FILL: Reference 2.]

\end{thebibliography}

% ============================================================================
% GLOSA AUDIT LAYER — INTERNAL OR SUPPLEMENTARY
% This layer protects epistemic traceability but is not forced into the visible
% manuscript unless the journal requests/permits it.
% ============================================================================

\ifglosaaudit
\appendix

\section{Glosa audit profile}
\label{app:glosa-profile}

\begin{auditblock}{A. Core epistemic roles}
\begin{description}[style=nextline,leftmargin=1.4em]
  \item[Experience-based / core respondent] [FILL: who brought the world/problem in]
  \item[Interactional expert / bridge] [FILL or None]
  \item[AI model(s)/tool(s)] [FILL: model + role]
  \item[Human accountable author/approver] [FILL]
\end{description}

\noindent Non-collapse:
\[
\text{Experience-Based Expertise}
\neq
\text{Interactional Expertise}
\neq
\text{AI Model}.
\]
\end{auditblock}

\begin{auditblock}{B. Standpoint and competence boundary}
\textbf{Standpoint:} [FILL]\\
\textbf{Disciplines/competencies explicitly not claimed:} [FILL]\\
\textbf{Local/practice knowledge scope:} [FILL]
\end{auditblock}

\begin{auditblock}{C. Question round trip}
\begin{description}[style=nextline,leftmargin=1.4em]
  \item[Question as lived] [FILL: verbatim human wording]
  \item[Question as readout] [FILL: what is actually observed/distinguished]
  \item[World-language hypothesis] [FILL]
  \item[Falsifier] [FILL: what record would defeat/narrow the claim]
\end{description}
\end{auditblock}

\section{Claim Matrix}
\label{app:claim-matrix}

\begin{longtable}{@{}p{0.05\linewidth}p{0.35\linewidth}p{0.20\linewidth}p{0.10\linewidth}p{0.10\linewidth}p{0.10\linewidth}@{}}
\toprule
ID & Claim & Evidence/formal support & Tier & Produced by & Responsible\\
\midrule
\endhead
C1 & [FILL] & [FILL] & [FILL] & human/AI/joint & human\\
C2 & [FILL] & [FILL] & [FILL] & human/AI/joint & human\\
\bottomrule
\end{longtable}

\section{Five-question claim audit}
\label{app:five-q}

For each load-bearing claim record:
\begin{enumerate}
  \item \texttt{seen}: what was actually observed/opened/executed;
  \item \texttt{separates}: what that record distinguishes;
  \item \texttt{ai\_filled}: what AI added;
  \item \texttt{assumed}: what was taken as given;
  \item \texttt{tested}: what independent check was run.
\end{enumerate}

\begin{longtable}{@{}p{0.05\linewidth}p{0.17\linewidth}p{0.17\linewidth}p{0.17\linewidth}p{0.17\linewidth}p{0.17\linewidth}@{}}
\toprule
ID & seen & separates & ai\_filled & assumed & tested\\
\midrule
\endhead
C1 & [FILL] & [FILL] & [FILL] & [FILL] & [FILL]\\
C2 & [FILL] & [FILL] & [FILL] & [FILL] & [FILL]\\
\bottomrule
\end{longtable}

\section{AI route-level disclosure and silent-lift audit}
\label{app:ai-route}

% Mirrors glosa P07 at a publication-compatible level.
\begin{longtable}{@{}p{0.26\linewidth}p{0.68\linewidth}@{}}
\toprule
Audit category & Record\\
\midrule
\endhead
Current evidence supplied/handled by AI & [FILL / none identified]\\
Retrieved tool evidence & [FILL / none identified]\\
Retained-record route & [FILL / none identified]\\
Model/calibration assumption & [FILL / none identified]\\
Prompt/system constraint & [FILL / none identified]\\
Decision policy proposed by AI & [FILL / none identified]\\
Represented dependency set & [FILL]\\
Actual dependency set & [FILL]\\
Silent-lift flags & [FILL: none / list]\\
Human completeness judgment & [FILL]\\
\bottomrule
\end{longtable}

\section{Literature Review System record}
\label{app:lrs}

% Compatible with glosa P13:
% L1 framing -> L2 frozen protocol -> L3 acquisition -> L4 extraction
% -> L5 citation verification -> L6 manifest/neighbour table.
\begin{tabularx}{\linewidth}{@{}lXl@{}}
\toprule
Stage & Required artifact & Status\\
\midrule
L1 & Frozen question/scope & [FILL]\\
L2 & Search protocol / honest review mode & [FILL]\\
L3 & Acquisition log & [FILL]\\
L4 & Source reading/extraction records & [FILL]\\
L5 & Metadata + claim-match verification & [FILL]\\
L6 & Manifest + neighbour/dialogue table & [FILL]\\
\bottomrule
\end{tabularx}

\medskip
\noindent\textbf{Accuracy gate:} [FILL: PASS / PASS WITH LIMITS / FAIL]\\
\textbf{Diversity gate:} [FILL: PASS / PASS WITH LIMITS / FAIL]\\
\textbf{Secondary-citation audit:} [FILL]\\
\textbf{Retraction/correction check:} [FILL]\\
\textbf{Review-mode claim:} [FILL; never call targeted search ``systematic'']

\section{Blackbox Note and transformation log}
\label{app:blackbox}

\begin{longtable}{@{}p{0.12\linewidth}p{0.12\linewidth}p{0.12\linewidth}p{0.42\linewidth}p{0.16\linewidth}@{}}
\toprule
ID & Role & Kind & Verbatim retained human line & Became\\
\midrule
\endhead
BB-001 & [FILL] & [FILL] & [FILL: exact retained wording] & [FILL: claim/rule]\\
\bottomrule
\end{longtable}

\medskip
\begin{longtable}{@{}p{0.12\linewidth}p{0.12\linewidth}p{0.18\linewidth}p{0.50\linewidth}@{}}
\toprule
Step & By & Input & Transformation / output\\
\midrule
\endhead
[FILL] & human/AI/joint & [FILL] & [FILL]\\
\bottomrule
\end{longtable}

\section{Knowledge state and independence}
\label{app:kstate}

\noindent\textbf{Overall K-state:} [FILL]\\
\textbf{Highest claim tier:} [FILL]\\
\textbf{Independence class(es):} [FILL]\\
\textbf{Known route dependence:} [FILL]\\
\textbf{Claims on HOLD/Open:} [FILL]

\section{Human Mastery Gate}
\label{app:mastery}

Before submission/public release, the human accountable author should be able
to defend without AI assistance:
\begin{enumerate}
  \item the problem;
  \item major literature families;
  \item closest constructs/frameworks;
  \item the precise inadequacy;
  \item the proposed mechanism/argument;
  \item boundary conditions;
  \item strongest counterargument;
  \item strongest counterexample;
  \item falsifying record;
  \item discriminating empirical/formal test where relevant.
\end{enumerate}

\noindent\textbf{Gate status:} [FILL: PASS / PASS WITH NAMED GAPS / NOT READY]\\
\textbf{Named gaps:} [FILL]

\section{Pre-publication adversarial gate}
\label{app:publish-gate}

% Based on glosa P10; upgraded to include publication-ethics declarations.
\begin{tabularx}{\linewidth}{@{}p{0.08\linewidth}p{0.63\linewidth}p{0.20\linewidth}@{}}
\toprule
Gate & Check & Status\\
\midrule
R1 & Leak/privacy/confidentiality scan & [FILL]\\
R2 & License, copyright, permissions, and IP coverage & [FILL]\\
R3 & Claim-tier / evidence-strength fidelity & [FILL]\\
R4 & Citation accuracy and source-match recheck & [FILL]\\
R5 & Lineage / anchor-preservation audit & [FILL]\\
R6 & Overclaim/register scan (contaminated-concept table) & [FILL]\\
R7 & Completeness / dangling-reference / supplement check & [FILL]\\
R8 & Human Mastery Gate & [FILL]\\
R9 & Authorship + CRediT + acknowledgements consistency & [FILL]\\
R10 & Funding + competing-interest disclosure completeness & [FILL]\\
R11 & Ethics/consent/data/code/protocol declarations & [FILL]\\
R12 & AI public disclosure matches internal AI-use log & [FILL]\\
R13 & Target-journal current policy checked & [FILL]\\
R14 & Confidential peer-review material not exposed to non-permitted AI & [FILL]\\
\bottomrule
\end{tabularx}

\section{Reproduction ledger}
\label{app:reproduction}

\begin{longtable}{@{}p{0.24\linewidth}p{0.70\linewidth}@{}}
\toprule
Object & Immutable/reproducible locator\\
\midrule
\endhead
Manuscript source & [FILL: commit/tag/hash]\\
Bibliography/source manifest & [FILL]\\
Data & [FILL / N/A]\\
Code & [FILL / N/A]\\
Environment & [FILL / N/A]\\
Formal proof & [FILL / N/A]\\
Figures/tables generation & [FILL / N/A]\\
AI prompt/output log & [FILL / restricted + reason]\\
Validation records & [FILL]\\
Zenodo/OSF/other DOI & [FILL / N/A]\\
\bottomrule
\end{longtable}

\section{Internal policy-interoperability checklist}
\label{app:policy-map}

This internal table is a routing aid, not a legal opinion and not a claim that
publishers have identical rules.

\begin{tabularx}{\linewidth}{@{}p{0.32\linewidth}XXXX@{}}
\toprule
Item & NRCT-TH & ICMJE & Elsevier & Springer/Wiley\\
\midrule
Human accountability & check & check & check & check\\
AI not author & check & check & check & check\\
Substantive AI disclosure & check & check & check & check\\
Research-process AI in methods & check & check & check & check\\
Citation/source verification & check & check & check & check\\
Privacy/confidentiality & check & check & check & check\\
IP/rights/terms review & check & -- & check & check\\
Prompt/output/validation log & check & recommended & recommended & documented\\
AI peer-review confidentiality & check & check & check & check\\
Competing-interest statement & journal-dependent & check & journal-dependent & check\\
Funding statement & journal-dependent & check & journal-dependent & check\\
\bottomrule
\end{tabularx}

\fi % end glosaaudit

% ============================================================================
% END-OF-FILE IMPLEMENTATION NOTES
% ============================================================================

\ifsubmissionmode
% Submission mode:
% - remove/internalize glosa-only metadata if the journal does not request it;
% - keep the seven intellectual functions even if headings are merged;
% - keep declarations required by the target journal;
% - re-check the current author guidelines on the actual submission date.
\else
% Internal mode:
% - retain all audit appendices;
% - retain model/version/time, prompt/output records subject to privacy/IP;
% - do not publicly expose confidential/sensitive prompts merely for
%   reproducibility.
\fi

\end{document}
